% =====================================================================
%  1bat-ud3-2.tex  —  HORA 2: Gràfiques mudes: llegir fenòmens socials
%  (sense preàmbul: s'inclou des de main.tex)
% =====================================================================
\horatitol{Sessió 2 — Gràfiques mudes: llegir fenòmens socials}%
{Aparellar cada gràfica amb el fenomen que representa i fer emergir el vocabulari de l'anàlisi de funcions.}

\subsection{Idea clau}

A la sessió anterior vam veure que una funció és una magnitud que \textbf{depèn}
d'una altra, i que es pot llegir amb una frase, una taula o una \textbf{gràfica}.
Avui ens centrem en la gràfica: la \emph{forma} d'una corba ja explica moltes coses
del fenomen, encara sense números. Una gràfica «muda» (sense valors als eixos) ens
permet preguntar-nos si una magnitud \textbf{creix} o \textbf{decreix}, si arriba a
un \textbf{màxim}, o si s'acosta a un límit sense tocar-lo.

\begin{idea}
\textbf{Vocabulari que farem servir per descriure una gràfica}
\begin{itemize}[leftmargin=1.4em, itemsep=2pt]
  \item \textbf{Creix}: a mesura que avancem cap a la dreta, la corba puja.
  \item \textbf{Decreix}: a mesura que avancem cap a la dreta, la corba baixa.
  \item \textbf{Màxim} (o mínim): un punt més alt (o més baix) que els del voltant.
  \item \textbf{S'acosta a un límit}: la corba s'aproxima a un valor sense arribar-hi
        mai (més endavant en direm \emph{asímptota}).
\end{itemize}
La regla d'or de la lectura social: \textbf{la forma ha de tenir sentit}. Un cost no
pot ser negatiu; un nombre de persones no pot baixar de zero.
\end{idea}

\subsection{Exemples}

\begin{exemple}[quatre formes per a quatre fenòmens]
Aquestes són quatre gràfiques mudes típiques. Fixa't només en la \emph{forma}:

\begin{center}
\begin{tikzpicture}
\begin{axis}[udaxis, width=6.6cm, height=4.6cm, xtick=\empty, ytick=\empty,
    xmin=0, xmax=10, ymin=0, ymax=10, title={\small\color{udnavy}(A)}]
  \addplot[udblue, domain=0:9.5]{0.9*x+0.5};
\end{axis}
\begin{axis}[udaxis, width=6.6cm, height=4.6cm, xtick=\empty, ytick=\empty,
    xmin=0, xmax=10, ymin=0, ymax=10, xshift=6.4cm, title={\small\color{udnavy}(B)}]
  \addplot[udblue, domain=0:10, samples=80]{-0.34*(x-5)^2+9};
\end{axis}
\end{tikzpicture}

\begin{tikzpicture}
\begin{axis}[udaxis, width=6.6cm, height=4.6cm, xtick=\empty, ytick=\empty,
    xmin=0, xmax=10, ymin=0, ymax=10, title={\small\color{udnavy}(C)}]
  \addplot[udblue, domain=0.6:10, samples=100]{5/x+0.3};
\end{axis}
\begin{axis}[udaxis, width=6.6cm, height=4.6cm, xtick=\empty, ytick=\empty,
    xmin=0, xmax=10, ymin=0, ymax=10, xshift=6.4cm, title={\small\color{udnavy}(D)}]
  \addplot[udblue, domain=0:10, samples=120]{5+2.6*sin(deg(x))};
\end{axis}
\end{tikzpicture}
\figcap{Quatre formes bàsiques: recta creixent, corba amb màxim, decreixent cap a un límit, i amb pujades i baixades.}
\end{center}

\noindent Llegim-les: \textbf{(A)} creix sempre al mateix ritme (una recta);
\textbf{(B)} creix, arriba a un \textbf{màxim} i després decreix; \textbf{(C)}
decreix de pressa i després s'acosta a un límit sense arribar-hi; \textbf{(D)}
puja i baixa diverses vegades.
\end{exemple}

\begin{exemple}[de la forma al fenomen]
Quin fenomen social podria correspondre a cada forma de l'exemple anterior?
\begin{itemize}[leftmargin=1.4em, itemsep=2pt]
  \item \textbf{(A)} El cost de la llum segons els kWh consumits: com més consum,
        més cost, i sempre al mateix preu per kWh $\rightarrow$ \emph{recta creixent}.
  \item \textbf{(C)} Un fons repartit entre cada cop més famílies: com més famílies,
        menys rep cadascuna, acostant-se a zero $\rightarrow$ \emph{decreixent cap a un límit}.
\end{itemize}
La mateixa forma matemàtica pot descriure realitats molt diferents; per això cal
\textbf{argumentar} per què una forma encaixa amb un fenomen.
\end{exemple}

\subsection{Exercicis resolts}

\begin{exresolt}[aparellar gràfica i fenomen]
Tenim tres fenòmens i tres gràfiques mudes. Aparella'ls i justifica:
\begin{enumerate}[label=\arabic*., leftmargin=1.6em]
  \item L'estalvi acumulat en una guardiola si hi poso la mateixa quantitat cada setmana.
  \item La temperatura d'un cafè que es va refredant fins a la temperatura de l'habitació.
  \item El nombre de contagiats d'un brot que creix, fa cim i després baixa.
\end{enumerate}

\begin{center}
\begin{tikzpicture}
\begin{axis}[udaxis, width=5.2cm, height=4cm, xtick=\empty, ytick=\empty,
    xmin=0, xmax=10, ymin=0, ymax=10, title={\small\color{udnavy}(I)}]
  \addplot[udblue, domain=0:9.5]{0.95*x+0.3};
\end{axis}
\begin{axis}[udaxis, width=5.2cm, height=4cm, xtick=\empty, ytick=\empty,
    xmin=0, xmax=10, ymin=0, ymax=10, xshift=5.0cm, title={\small\color{udnavy}(II)}]
  \addplot[udblue, domain=0:10, samples=100]{2+7*exp(-0.5*x)};
\end{axis}
\begin{axis}[udaxis, width=5.2cm, height=4cm, xtick=\empty, ytick=\empty,
    xmin=0, xmax=10, ymin=0, ymax=10, xshift=10.0cm, title={\small\color{udnavy}(III)}]
  \addplot[udblue, domain=0:10, samples=80]{-0.34*(x-5)^2+9};
\end{axis}
\end{tikzpicture}
\figcap{Tres gràfiques mudes per aparellar.}
\end{center}

\solucio
Mirem quina \emph{forma} té sentit per a cada història:
\begin{itemize}[leftmargin=1.4em, itemsep=2pt]
  \item L'estalvi puja sempre al mateix ritme (mateixa quantitat cada setmana):
        forma de \textbf{recta creixent} $\rightarrow$ \textbf{(I)}.
  \item El cafè baixa de pressa al principi i s'acosta a la temperatura de
        l'habitació sense baixar-ne més: \textbf{decreix cap a un límit}
        $\rightarrow$ \textbf{(II)}.
  \item El brot puja, fa \textbf{màxim} i baixa: corba amb un cim $\rightarrow$ \textbf{(III)}.
\end{itemize}
\resposta{$1\rightarrow$ (I), \quad $2\rightarrow$ (II), \quad $3\rightarrow$ (III).}
\end{exresolt}

\begin{exresolt}[descriure amb el vocabulari correcte]
Observa la gràfica del repartiment d'un fons d'emergència entre $x$ famílies.
Descriu-la amb el vocabulari de la idea clau i interpreta-la socialment.

\begin{center}
\begin{tikzpicture}
\begin{axis}[udaxis, width=8cm, height=5cm, xtick=\empty, ytick=\empty,
    xlabel={\small nombre de famílies}, ylabel={\small \euro\ per família},
    xmin=0, xmax=10, ymin=0, ymax=10]
  \addplot[udblue, domain=0.55:10, samples=120]{5/x};
  \draw[udline, dashed] (axis cs:0,0) -- (axis cs:10,0);
\end{axis}
\end{tikzpicture}
\figcap{Diners que rep cada família segons el nombre de famílies.}
\end{center}

\solucio
La corba \textbf{decreix} sempre: com més famílies, menys rep cadascuna. Al
principi baixa molt de pressa i després cada cop més a poc a poc, \textbf{acostant-se
a zero} sense arribar-hi mai. Socialment: si el fons s'ha de repartir entre moltes
famílies, l'ajut per família es fa tan petit que pràcticament desapareix; els
recursos «s'esgoten» en repartir-los massa.
\resposta{Funció \emph{decreixent} que s'acosta a zero sense tocar-lo; com més
famílies, menys ajut per cadascuna.}
\end{exresolt}

\subsection{Exercicis per fer}

\begin{exercicis}
\begin{exlist}
  \item Per a cada gràfica, digues si \textbf{creix}, \textbf{decreix} o té un
        \textbf{màxim}:

        \begin{center}
        \begin{tikzpicture}
        \begin{axis}[udaxis, width=5cm, height=3.8cm, xtick=\empty, ytick=\empty,
            xmin=0, xmax=10, ymin=0, ymax=10, title={\small\color{udnavy}(a)}]
          \addplot[udblue, domain=0:10, samples=80]{-0.32*(x-6)^2+9.5};
        \end{axis}
        \begin{axis}[udaxis, width=5cm, height=3.8cm, xtick=\empty, ytick=\empty,
            xmin=0, xmax=10, ymin=0, ymax=10, xshift=4.8cm, title={\small\color{udnavy}(b)}]
          \addplot[udblue, domain=0:9.5]{-0.85*x+9};
        \end{axis}
        \begin{axis}[udaxis, width=5cm, height=3.8cm, xtick=\empty, ytick=\empty,
            xmin=0, xmax=10, ymin=0, ymax=10, xshift=9.6cm, title={\small\color{udnavy}(c)}]
          \addplot[udblue, domain=0:9.5]{0.8*x+1};
        \end{axis}
        \end{tikzpicture}
        \figcap{Tres gràfiques per descriure.}
\end{center}

  \item Dibuixa \emph{tu} (a mà, en la teva llibreta) una gràfica muda que
        representi cada situació, sense números:
        \begin{enumerate}[label=\alph*), leftmargin=1.6em, topsep=2pt]
          \item El cost total d'un casal segons el nombre d'infants (preu fix per infant).
          \item L'aigua que queda en un dipòsit que es va buidant.
        \end{enumerate}
  \item Associa cada frase amb la forma que li correspon (recta creixent,
        decreixent cap a un límit, o corba amb màxim):
        \begin{enumerate}[label=\alph*), leftmargin=1.6em, topsep=2pt]
          \item «Els diners recaptats en una rifa segons els bitllets venuts.»
          \item «El rendiment d'una activitat segons les hores: millora, fa cim i
                empitjora per cansament.»
          \item «La concentració d'un medicament al cos, que va baixant amb les hores.»
        \end{enumerate}
  \item Mira la gràfica de l'evolució de l'atur en una comarca durant uns anys.
        Indica \textbf{un} tram en què creix i \textbf{un} tram en què decreix.

        \begin{center}
        \begin{tikzpicture}
        \begin{axis}[udaxis, width=8.6cm, height=4.6cm, xtick=\empty, ytick=\empty,
            xlabel={\small anys}, ylabel={\small atur},
            xmin=0, xmax=10, ymin=0, ymax=10]
          \addplot[udblue, domain=0:10, samples=160]{5+2.6*sin(deg(0.9*x))-0.15*x+1};
        \end{axis}
        \end{tikzpicture}
        \figcap{Evolució de l'atur al llarg dels anys (gràfica muda).}
\end{center}

  \item \textbf{(Compte amb la lectura)} Algú diu: «aquesta gràfica baixa, per tant
        el fenomen és dolent». Explica per què aquesta afirmació pot ser falsa, posant
        un exemple social en què una gràfica \emph{decreixent} sigui una bona notícia.
  \item \textbf{(Raonament)} Tria un fenomen del teu entorn (àmbit social, casa,
        institut) i descriu-ne l'evolució amb el vocabulari d'avui: digues si creix,
        decreix, si té algun màxim o si s'acosta a algun límit. No cal cap número.
\end{exlist}
\end{exercicis}

% ---------------------------------------------------------------------
\fullsolucions{Sessió 2}
\begin{solucions}
\begin{sollist}
  \item (a) Té un \textbf{màxim} (puja i després baixa). \quad
        (b) \textbf{Decreix}. \quad (c) \textbf{Creix}.
  \item Resposta gràfica oberta. (a) Recta creixent que parteix de l'origen (cost
        proporcional als infants). (b) Corba \emph{decreixent} fins a zero (el dipòsit
        es buida). Es valora la forma coherent, no els números.
  \item (a) Recta creixent. \quad (b) Corba amb màxim. \quad
        (c) Decreixent cap a un límit (s'acosta a zero).
  \item Resposta oberta segons lectura: hi ha trams de pujada (creix) i de baixada
        (decreix) alternats; n'hi ha prou d'assenyalar-ne un de cada amb coherència.
  \item Resposta oberta. Exemple vàlid: la gràfica del nombre de persones sense llar
        ateses que encara no tenen habitatge \emph{decreix} si el servei funciona; o el
        temps d'espera d'una llista que baixa. Decréixer no és, en si mateix, dolent.
  \item Resposta oberta. Es valora l'ús correcte del vocabulari (creix/decreix/màxim/límit)
        i la coherència amb el fenomen triat.
\end{sollist}
\end{solucions}
